\section{Related Works}
\label{sec:related_works}

\subsection{Demographic Face Recognition}
\label{sec:demographic_face_recognition}

Perkembangan analisis atribut demografis wajah bermula dari pemanfaatan deskriptor geometri lokal serta fitur buatan tangan seperti Local Binary Patterns dan filter Gabor menuju representasi mendalam berbasis CNN. Pendekatan berbasis wilayah lokal memanfaatkan geometri area wajah tengah untuk mengenali etnisitas sekaligus menekan interferensi latar belakang citra~\cite{ref11}. Namun, arsitektur konvolusi klasik masih menghadapi tantangan ketahanan terhadap fluktuasi pencahayaan ekstrem dan disparitas performa antarras yang dikenal sebagai \textit{Other-Race Effect}. Peningkatan kedalaman jaringan terbukti mampu memperluas daya diskriminasi fitur visual dalam membedakan populasi heterogen~\cite{ref13}. Selain pendalaman lapisan, mekanisme berbagi fitur paralel dengan modul atensi visual dikembangkan untuk menangkap korelasi antaratribut wajah~\cite{ref16}, memperkuat transisi ke arah representasi yang lebih adaptif terhadap keragaman morfologi wajah.

Setelah ekstraksi fitur, penentuan pengklasifikasi hilir memegang peranan krusial dalam memetakan representasi laten ke ruang keputusan demografis. Model konvensional umumnya mengandalkan lapisan SoftMax secara \textit{end-to-end} untuk menghasilkan probabilitas kelas~\cite{ref11}, \cite{ref18}. Namun, pemetaan nonlinier seperti Kernel Extreme Learning Machine (KELM) dan SVM pada representasi beku menawarkan batas pemisahan kokoh tanpa pembaruan bobot \textit{backbone}~\cite{ref13}, \cite{ref12}. Formulasi matematis fungsi keputusan KELM dimodelkan pada \eqref{eq:1}:
\begin{equation}
f(\mathbf{x}) = \mathbf{h}(\mathbf{x})\boldsymbol{\beta} = \sum_{i=1}^M \beta_i K(\mathbf{x}, \mathbf{x}_i)
\label{eq:1}
\end{equation}
di mana $\mathbf{h}(\mathbf{x})$ menyatakan vektor pemetaan fitur laten, $\boldsymbol{\beta} = [\beta_1, \dots, \beta_M]^T$ melambangkan bobot teroptimasi, $K(\mathbf{x}, \mathbf{x}_i)$ merupakan fungsi kernel nonlinier terhadap sampel pendukung $\mathbf{x}_i$, dan $M$ menyatakan jumlah basis representasi. Pemetaan kernel ini memungkinkan pemisahan batas keputusan nonlinier yang efektif pada ruang berdimensi tinggi tanpa optimasi \textit{gradient descent} berulang.

\subsection{Vision Transformers in Face Analysis}
\label{sec:vision_transformers_face}

Penerapan ViT dalam biometrik wajah menghadirkan pemodelan global yang mengatasi bias induktif lokal CNN~\cite{ref14}. Citra masukan dipartisi menjadi $N$ \textit{patch} beresolusi $P \times P$ dengan $C$ kanal warna, lalu diproyeksikan ke dimensi laten via formulasi \eqref{eq:2}:
\begin{equation}
\mathbf{x}_p \in \mathbb{R}^{N \times (P^2 \cdot C)}, \quad \mathbf{z}_0 = [\mathbf{x}_{\text{class}}; \, \mathbf{x}_p^1\mathbf{E}; \dots; \, \mathbf{x}_p^N\mathbf{E}] + \mathbf{E}_{\text{pos}}
\label{eq:2}
\end{equation}
di mana $N = HW/P^2$ menyatakan jumlah \textit{patch}, $P$ ukuran \textit{patch}, $C$ kanal warna, $\mathbf{E} \in \mathbb{R}^{(P^2 \cdot C) \times D}$ matriks proyeksi, $\mathbf{x}_{\text{class}}$ token kelas, dan $\mathbf{E}_{\text{pos}}$ penyandian posisi. Interaksi spasial global antarpatch dimodelkan melalui mekanisme MHSA sesuai \eqref{eq:3}:
\begin{equation}
\text{Attention}(\mathbf{Q}, \mathbf{K}, \mathbf{V}) = \text{softmax}\left(\frac{\mathbf{Q}\mathbf{K}^T}{\sqrt{d_k}}\right)\mathbf{V}
\label{eq:3}
\end{equation}
dengan $\mathbf{Q}$, $\mathbf{K}$, dan $\mathbf{V}$ merepresentasikan matriks \textit{Query}, \textit{Key}, dan \textit{Value}, serta $\sqrt{d_k}$ faktor penskalaan. Atensi ini efektif diterapkan pada klasifikasi gender wajah~\cite{ref17}, sementara MaxViT memperkaya hierarki representasi lewat atensi blok dan kisi~\cite{ref18}.

Karakteristik representasi global yang dihasilkan oleh arsitektur transformer mendorong telaah mendalam terkait propagasi bias demografis dalam sistem visi komputer. Pengujian pada korpus data yang seimbang menyediakan kondisi eksperimental terkontrol untuk membandingkan disparitas performa antarsubkelompok tanpa distorsi ketimpangan jumlah sampel kelas. Walaupun distribusi data disetarakan, keseimbangan tersebut tidak secara otomatis meniadakan bias sosial yang melekat pada model. Investigasi representasi laten membuktikan bahwa arsitektur transformer dapat mereproduksi serta melipatgandakan bias sosial yang bersumber dari korpus pra-latih berskala masif~\cite{ref14}. Sebagai langkah penanganan ketimpangan tersebut, pendekatan pandangan generatif berbasis StyleGAN2 dimanfaatkan untuk memperkaya keragaman fenotipik antarkelompok demografis~\cite{ref15}. Integrasi metode ini membuktikan bahwa pengujian keadilan representasi dan penyelarasan distribusi data sangat krusial dalam perancangan model biometrik wajah modern.

\subsection{Multi-Domain Feature Fusion}
\label{sec:multidomain_feature_fusion}

Keterbatasan representasi domain tunggal mendorong perkembangan pembelajaran multi-tugas dan fusi lintas representasi. Pembelajaran multi-tugas terbukti efektif memetakan korelasi estimasi usia dan gender secara terpadu~\cite{ref12}, sementara skema berbagi representasi memfasilitasi ekstraksi simultan berbagai atribut afektif~\cite{ref16}. Selanjutnya, integrasi fitur laten berbasis ViT membuktikan keunggulan penggabungan domain biometrik dengan ekspresi emosi~\cite{ref19} serta usia biologis~\cite{ref20}. Fusi representasi multi-domain ortogonal tersebut diformulasikan melalui operasi direct sum pada \eqref{eq:4}:
\begin{equation}
\mathbf{z}_{\text{fused}} = \mathbf{z}_{\text{face}} \oplus \mathbf{z}_{\text{emotion}} \oplus \mathbf{z}_{\text{age}} \in \mathbb{R}^{D_{\text{total}}}
\label{eq:4}
\end{equation}
di mana operator $\oplus$ menyatakan konkatenasi langsung vektor fitur berbobot sama dari domain biometrik $\mathbf{z}_{\text{face}}$, ekspresi afektif $\mathbf{z}_{\text{emotion}}$, dan usia biologis $\mathbf{z}_{\text{age}}$, dengan dimensi gabungan $D_{\text{total}} = 2.304$. Integrasi ini memadukan isyarat identitas statis, variasi ekspresi dinamis, dan pola penuaan alami secara komprehensif.

Mempertimbangkan batasan metodologis pada studi terdahulu, penelitian ini menegaskan posisi kebaruan melalui fusi tri-domain terpadu yang memadukan representasi biometrik wajah, ekspresi emosi afektif, dan usia biologis dari model Vision Transformer pra-latih dengan bobot \textit{backbone} beku. Berbeda dari literatur sebelumnya yang mengevaluasi performa pada skema pembagian citra acak berisiko kebocoran identitas, kerangka kerja ini dievaluasi secara ketat di bawah protokol \textit{Strict Subject-Disjoint Partitioning} untuk menjamin keabsahan generalisasi pada subjek yang belum pernah dilihat sebelumnya. Selain itu, kontribusi penelitian ini mencakup eksplorasi sistematis tujuh skema ablasi representasi fitur lintas empat famili pengklasifikasi pembelajaran mesin klasik, optimasi prapemrosesan terisolasi tingkat lipatan bebas kebocoran, serta karakterisasi inferensial formal mengenai fenomena saturasi representasi nonlinier versus komplementaritas linier.
